\section{Introducción}

% -------------------------------------------------------------------------
% Diapositiva 1: Contexto y Motivación
% -------------------------------------------------------------------------
\begin{frame}{Contexto y Motivación}
    \begin{columns}
        \begin{column}{0.6\textwidth}
            \textbf{Motores BLDC y PMSM:}
            \begin{itemize}
                \item Alta eficiencia y densidad de potencia.
                \item Aplicaciones críticas: Vehículos eléctricos, robótica, aeroespacial.
            \end{itemize}
            
            \vspace{0.5cm}
            
            \textbf{El Problema del Rizado de Par (Torque Ripple):}
            \begin{itemize}
                \item Vibraciones mecánicas y ruido acústico.
                \item Desgaste prematuro de componentes.
                \item \alert{Causa principal:} Interacción entre corrientes no ideales y la fuerza contraelectromotriz (Back-EMF) con armónicos.
            \end{itemize}
        \end{column}
        \begin{column}{0.4\textwidth}
            \begin{figure}
                \centering
                % SUGERENCIA: Imagen genérica de un BLDC o diagrama de aplicación
                \includegraphics[width=\textwidth]{images/bldc_application_example.png}
                \caption{Aplicaciones de alto rendimiento.}
            \end{figure}
        \end{column}
    \end{columns}
\end{frame}

% -------------------------------------------------------------------------
% Diapositiva 2: Estrategias de Control
% -------------------------------------------------------------------------
\begin{frame}{Evolución de las Estrategias de Control}
    \begin{block}{Control Clásico (PI - FOC)}
        \begin{itemize}
            \item Modulador PWM fijo.
            \item Limitado por el ancho de banda del lazo lineal.
            \item Dificultad para manejar restricciones y no-linealidades.
        \end{itemize}
    \end{block}

    \vfill

    \begin{block}{Control Predictivo (FCS-MPC) \hfill \textit{Enfoque de esta Tesis}}
        \begin{itemize}
            \item Aprovecha la naturaleza discreta del inversor (Estados finitos).
            \item \textbf{Ventaja:} Respuesta dinámica superior y manejo intuitivo de restricciones.
            \item \textbf{Desafío:} \alert{Alta dependencia del modelo matemático.}
        \end{itemize}
    \end{block}
\end{frame}

% -------------------------------------------------------------------------
% Diapositiva 3: El Desafío del "Model Mismatch"
% -------------------------------------------------------------------------
\begin{frame}{El Desafío: Discrepancia del Modelo (Mismatch)}
    El rendimiento del FCS-MPC depende de la precisión de la predicción:
    
    \begin{equation*}
        i(k+1) = i(k) + \frac{T_s}{L} \left( u(k) - \alert{e(k)} - R i(k) \right)
    \end{equation*}

    \begin{columns}
        \begin{column}{0.5\textwidth}
            \textbf{Modelo Ideal (Controlador):}
            $$ \hat{e}(k) \approx k_e \omega \sin(\theta_e) $$
            \textit{"El controlador asume un mundo perfecto."}
        \end{column}
        \begin{column}{0.5\textwidth}
            \textbf{Planta Real (Motor Físico):}
            $$ e(k) = \sum_{n=1,5,7...} A_n \sin(n\theta_e + \phi_n) + \text{Ruido} $$
            \textit{"La realidad tiene armónicos y perturbaciones."}
        \end{column}
    \end{columns}

    \vspace{0.5cm}
    \centering
    \small \textbf{Consecuencia:} Si $\hat{e}(k) \neq e(k)$, surge un error de predicción que se traduce en rizado de corriente.
\end{frame}

% -------------------------------------------------------------------------
% Diapositiva 4: Análisis Preliminar de Datos (Tu aporte con Python)
% -------------------------------------------------------------------------
\begin{frame}{Evidencia Experimental: Caracterización de BEMF}
    Se analizaron datos experimentales de una prueba de arrastre (Open Circuit) para identificar la "huella" del motor.
    
    \begin{columns}
        \begin{column}{0.5\textwidth}
            \begin{itemize}
                \item \textbf{Método:} FFT sobre datos normalizados por velocidad.
                \item \textbf{Hallazgo:} Presencia significativa de armónicos $5^\circ$ ($2.2\%$) y $7^\circ$ ($1.1\%$).
                \item Estos armónicos generan oscilaciones de par de orden $6\omega_e$.
            \end{itemize}
        \end{column}
        \begin{column}{0.5\textwidth}
            \begin{figure}
                \centering
                % AQUÍ VA LA GRÁFICA QUE HICIMOS EN PYTHON (bemf_final.png)
                % La que muestra la forma de onda reconstruida y las barras de FFT
                \includegraphics[width=\textwidth]{images/bemf_characterization_python.png}
                \caption{Análisis espectral de la Back-EMF real.}
            \end{figure}
        \end{column}
    \end{columns}
\end{frame}

% -------------------------------------------------------------------------
% Diapositiva 5: Objetivos y Propuesta
% -------------------------------------------------------------------------
\begin{frame}{Propuesta y Objetivos de la Tesis}
    \textbf{Objetivo General:}
    Diseñar y validar estrategias de FCS-MPC robustas ante incertidumbres de la Back-EMF para minimizar el rizado de corriente.

    \vspace{0.5cm}
    
    \textbf{Metodología (Gemelo Digital):}
    \begin{enumerate}
        \item \textbf{FCS-M2PC Mejorado (Offline):} Inyección de armónicos $H_5$ y $H_7$ en el modelo de predicción (Basado en datos).
        \item \textbf{Observador de Modos Deslizantes (Online):} Implementación de un diferenciador de \textit{Levant (Super-Twisting)} para estimar la BEMF desconocida en tiempo real.
        \item \textbf{Validación:} Comparación en simulación de alta fidelidad incluyendo ruido de medición y efectos parásitos.
    \end{enumerate}
\end{frame}